\subsection{Graphing Quadratic Functions in Standard Form} 
\begin{definition}
A quadratic function written as $$f(x)=ax^2+bx+c$$ where $a\neq 0$, $b$, and $c$ are real numbers is said to be written in \term{standard form}.
\end{definition}
To graph a quadratic function given in standard form, we will use the following procedure.
\begin{procedure}{Graphing a Quadratic Function in Standard Form}
\begin{enumerate*}
\item Create a table of values showing (integer) pairs $x$ and $y=f(x)$. Include enough $x$ values that $y$ values both increase and decrease.
\item If possible, find two pairs with the same $y$-coordinate. The axis of symmetry lies halfway between their $x$-coordinates.
\item Find the coordinates of the vertex, if it doesn't have integer coordinates.
\item Plot the vertex, along with at least two points on each side of the axis of symmetry.
\item Draw a parabola passing through the two plotted points.
\end{enumerate*}
\end{procedure}
This process will require a certain amount of trial and error to find ``appropriate'' points! (We will have a more efficient method once we know more about quadratic functions.)\
\begin{Example}
Fill in the table below. Then, sketch the graph of $f(x)=x^2-4x+3$.
\begin{lpic}{}{10}
\GraphIn{4}{}
\regtext{1}{-1}{$x$};
\regtextleft{3}{-1}{$y=f(x)$};
\draw (0,-1) -- (6,-1);
\draw (2,0) -- (2,-10);
\foreach \x in {-4,...,4} {
	\regtext{1}{-6-\x}{$\x$};
}
\regtextleft{7}{-2}{Vertex:};
\regtextleft{7}{-1}{AoS:};
\end{lpic}
\end{Example}